\chapter{示例设备与工作流程}\label{ch:device}
\index{示例设备}

本章用一个虚构的设备说明技术书里最常见的“结构、图、表、公式和单位”如何排在一起。所有数据只用于排版练习，不代表真实设备。

\section{设备组成}

示例设备由控制器、传感器、执行器和记录单元组成，工作流程见\cref{fig:process}。图中不写真实型号，只保留关系。

\begin{figure}[htbp]
  \centering
  \begin{tikzpicture}[
  node distance=1.1cm and 1.2cm,
  box/.style={draw=bookteal, fill=booksoft, rounded corners=2pt, minimum width=2.7cm, minimum height=1.1cm, align=center, font=\small},
  arrow/.style={-{Latex[length=2.4mm]}, thick, draw=bookteal}
]
  \node[box] (input) {采集输入};
  \node[box, right=of input] (check) {检查条件};
  \node[box, right=of check] (run) {执行任务};
  \node[box, below=of run] (record) {记录结果};
  \draw[arrow] (input) -- (check);
  \draw[arrow] (check) -- (run);
  \draw[arrow] (run) -- (record);
  \draw[arrow] (record.west) -- ++(-5.3,0) |- (input.south);
\end{tikzpicture}

  \caption{示例设备的任务流程}\label{fig:process}
\end{figure}

\section{单位与数值}

示例杆件长度记为 \qty{12.5}{\milli\metre}，环境温度记为 \qty{23}{\celsius}，热流密度记为 \qty{4.2}{\kilo\watt\per\square\metre}。表格中的数值用 \code{siunitx} 的 S 列对齐，见\cref{tab:si}。

\begin{table}[htbp]
  \centering
  \caption{示例测量值与单位}\label{tab:si}
  \begin{tabular}{@{}S[table-format=2.1]S[table-format=1.1]l@{}}
  \toprule
  {长度 / \si{\milli\metre}} & {温度 / \si{\celsius}} & 备注 \\
  \midrule
  12.5 & 23.0 & 示例值 \\
  8.0  & 22.4 & 示例值 \\
  15.2 & 24.1 & 示例值 \\
  \bottomrule
\end{tabular}

\end{table}

\section{公式与定义}

示例关系用\cref{eq:demo}表示：
\begin{equation}\label{eq:demo}
  q = \frac{P}{A},
\end{equation}
其中 $q$ 是热流密度，$P$ 是功率，$A$ 是面积。

\begin{definition}[示例杆件]
  示例杆件是一根只承受轴向载荷的虚构杆件，长度记为 $L$。
\end{definition}

\begin{example}
  若 $P=\qty{50}{\watt}$、$A=\qty{0.02}{\square\metre}$，则 $q=\qty{2500}{\watt\per\square\metre}$。计算过程要保留单位，不能只写数字。
\end{example}

\section{提示框与代码}

\begin{tcolorbox}[bookbox,title=课堂提示]
  正式项目里，单位、数字和技术结论必须回到原始 Word 或测试记录核对；排版改善不等于内容改写。\index{核对}
\end{tcolorbox}

下面的代码只演示等宽排版和行号，不包含真实项目逻辑：

\begin{lstlisting}[language=Python,caption={示例代码：只演示排版},label={lst:demo}]
def check_reading(value, limit):
    """Return True when a fictional reading is inside the limit."""
    return abs(value) <= limit

reading = 12.5
print(check_reading(reading, limit=15.0))
\end{lstlisting}


\section{本章小结}

本章把TikZ 流程图、\code{siunitx} 单位、编号公式、定义、示例、提示框和代码放在同一章里。下一章再补安装、子图、横向页和长表。\index{流程图}
